% Source printed pp. 293-301 / PDF pp. 301-309.
\section{依 Kollár 方法证明射影性}\label{src:6.9}
\begin{epigraph}
然而, 你也许对熟悉的旧有簇仍怀有感情. 尤其看来, \underline{射影}簇在代数几何中发挥着如此核心的技术作用, 即使你想摆脱它们, 实际上也可能办不到. 无论如何, 在可能时证明 [这个空间] 是射影簇, 是非常有趣的.
\epigraphauthor{David Mumford~\cite[441 页]{Mum76}}
\end{epigraph}
遵循 Kollár 在~\cite{Kol90} 中提出的方法, 我们给出 $\overline M_{g,\allowbreak n}$ 射影性的第二个证明; 这一方法建立在 Viehweg 的早期思想~\cite{Vie95} 之上. 为介绍证明射影性的总体策略, 先引入一些术语. 设 $\pi:\mathcal U_g\to\overline{\mathcal M}_g$ 为 $\operatorname{Spec}\mathbb Z$ 上的万有族. 对每个整数 $k\geq1$, 定义 $\overline{\mathcal M}_g$ 上的第 $k$ 个多重典范丛为
\[
\pi_*(\omega_{\mathcal U_g/\overline{\mathcal M}_g}^{\otimes k}).
\]
由命题~\bookref{6.3.23}, 它是向量丛, 当 $k=1$ 时秩为 $g$, 当 $k>1$ 时秩为 $(2k-1)(g-1)$. 其行列式
\[
\lambda_k:=\det\pi_*(\omega_{\mathcal U_g/\overline{\mathcal M}_g}^{\otimes k})
\]
是 $\overline{\mathcal M}_g$ 上的线丛. 当 $k=1$ 时, 它是定义~\bookref{6.10.1} 中的 Hodge 线丛.

\noindent\textbf{证明射影性的策略:} 证明当 $k\gg0$ 时, $\lambda_k$ 的某个正幂下降为粗模空间 $\overline M_g$ 上的丰沛线丛.
\subsection{Kollár 判据与稳定曲线粗模空间的射影性}\label{src:6.9.1-subsection}
Kollár 的策略是利用多重典范丛上的乘法映射
\begin{equation}\label{src:6.9.1}
W:=\operatorname{Sym}^m\underbrace{\pi_*(\omega_{\mathcal U_g/\overline{\mathcal M}_g}^{\otimes k})}_{V}
\longrightarrow\pi_*(\omega_{\mathcal U_g/\overline{\mathcal M}_g}^{\otimes mk})=:Q.
\end{equation}
它依赖与 Hilbert 概形的 GIT 构造完全相同的两个整数 $k,\allowbreak m$. 在稳定曲线 $C\in\mathcal M_g(k)$ 上, 此乘法映射的纤维为
\[
\operatorname{Sym}^m\mathrm H^0(C,\omega_C^{\otimes k})\longrightarrow\mathrm H^0(C,\omega_C^{\otimes mk}),
\]
其核由切出 $|\omega_C^{\otimes k}|:C\to\mathbb P(\mathrm H^0(C,\allowbreak \omega_C^{\otimes k}))$ 的像的 $m$ 次方程组成. 若 $k\geq3$, 则 $\omega_C^{\otimes k}$ 甚丰沛, 且当 $m\gg0$ 时, 可从乘法映射~\bookeqref{6.9.1} 的核恢复 $C$, 因而该乘法映射满射. 令 $w=\operatorname{rk}W$, $q=\operatorname{rk}Q$, $v=\operatorname{rk}V$. 因为 $W$ 表为对称积 $\operatorname{Sym}^mV$, 它具有沿 $\operatorname{GL}_v\to\operatorname{GL}_w$ 的结构群约化 (见定义~\bookref{B.1.64}).

对稳定曲线 $C$, 固定一个基 $\mathrm H^0(C,\allowbreak \omega_C^{\otimes k})\cong k^v$, 得到商
\[
(\operatorname{Sym}^mk^v\twoheadrightarrow\mathrm H^0(C,\omega_C^{\otimes mk}))\in\operatorname{Gr}(q,\operatorname{Sym}^mk^v).
\]
改变基的选择, 只会使这个商相差一个 $\operatorname{PGL}_v$ 作用. 此构造推广到稳定曲线族, 定义一个代数叠态射, 称为\newterm{分类映射}{classifying map}:
\[
\begin{aligned}
\overline{\mathcal M}_g&\longrightarrow[\operatorname{Gr}(q,\operatorname{Sym}^mk^v)/\operatorname{PGL}_v],\\
[C]&\longmapsto[\operatorname{Sym}^mk^v\cong\operatorname{Sym}^m\mathrm H^0(C,\omega_C^{\otimes k})\longrightarrow\mathrm H^0(C,\omega_C^{\otimes mk})].
\end{aligned}
\]
核心思想是借助 $\operatorname{Gr}(q,\allowbreak \operatorname{Sym}^mk^v)$ 的射影性证明 $\overline M_g$ 射影. 为此需要两个性质: 分类映射拟有限, 以及多重典范丛 $V=\pi_*(\omega_{\mathcal U_g/\overline{\mathcal M}_g}^{\otimes k})$ 数值有效.
\begin{theorem}[Kollár 丰沛性判据]\label{src:6.9.2}
设 $X$ 为域 $k$ 上的固有代数空间. 设 $W=\operatorname{Sym}^mV\twoheadrightarrow Q$ 为向量丛满射, 其中 $w=\operatorname{rk}W$, $q=\operatorname{rk}Q$, $v=\operatorname{rk}V$. 假设
\begin{enumerate}[(a)]
\item 分类映射
\[
\begin{aligned}
X&\longrightarrow[\operatorname{Gr}(q,\operatorname{Sym}^mk^v)/\operatorname{PGL}_v],\\
x&\longmapsto[W\otimes\kappa(x)\twoheadrightarrow Q\otimes\kappa(x)]
\end{aligned}
\]
拟有限;
\item $V$ 数值有效.
\end{enumerate}
则 $\det Q$ 丰沛.
\end{theorem}
我们将在~\bookref{6.9.3-subsection} 节详述这些假设并证明判据, 但先说明它如何推出 $\overline M_{g,\allowbreak n}$ 的射影性.
\begin{theorem}[\texorpdfstring{$\overline M_g$}{稳定曲线粗模空间}的射影性]\label{src:6.9.3}
当 $g\geq2$ 且 $N$ 充分可除时, $\overline{\mathcal M}_g$ 上的线丛 $\lambda_{18g-18}^{\otimes N}$ 下降为 $\overline M_g$ 上的线丛, 相对于 $\operatorname{Spec}\mathbb Z$ 丰沛. 特别地, $\overline M_g$ 在 $\operatorname{Spec}\mathbb Z$ 上射影.
\end{theorem}
\begin{proof}
由稳定曲线模叠的固有性 (\bookref{6.5.23}), $\overline{\mathcal M}_g$ 是 $\operatorname{Spec}\mathbb Z$ 上的固有 Deligne--Mumford 叠, 粗模空间 $\overline M_g$ 是 $\operatorname{Spec}\mathbb Z$ 上的固有代数空间. 由习题~\bookref{5.4.29}, 当 $N$ 充分可除时, $\lambda_{18g-18}^{\otimes N}$ 下降为 $\overline M_g$ 上的线丛 $L$. 因丰沛性是开性质 (\bookref{5.6.17}), 只需证明 $\overline M_g$ 定义于域 $k$ 上的情形. 由命题~\bookref{6.3.23}, $\omega_{\mathcal U_g/\overline{\mathcal M}_g}^{\otimes3}$ 相对甚丰沛, 且其推前是秩为 $5g-5$ 的向量丛. 对亏格 $g$ 的稳定曲线 $C$, 完全线性系 $|\omega_C^{\otimes3}|$ 将 $C$ 嵌入为 $\mathbb P^{5g-6}$ 中次数 $m:=6g-6$ 的曲线. 因为 $C$ 由 $m$ 次方程切出, 乘法映射
\[
W:=\operatorname{Sym}^m\pi_*(\omega_{\mathcal U_g/\overline{\mathcal M}_g}^{\otimes3})\twoheadrightarrow\pi_*(\omega_{\mathcal U_g/\overline{\mathcal M}_g}^{\otimes3m})=:Q
\]
是 $\overline{\mathcal M}_g$ 上的向量丛满射, $W$ 的秩为 $w:=\binom{5g-6+m}{m}$, $Q$ 的秩为 $q:=(6m-1)(g-1)$.

分类映射 $\overline{\mathcal M}_g\to[\operatorname{Gr}(q,\allowbreak k^w)/\operatorname{PGL}_{5g-5}]$ 在 $k$-点上单射, 因乘法映射的核唯一确定稳定曲线. 又因稳定曲线的自同构群 $\operatorname{Aut}(C)$ 等同于 $\operatorname{PGL}_{5g-5}$ 作用下 $[C\subseteq\mathbb P^{5g-6}]\in\operatorname{Hilb}^P(\mathbb P^{5g-6})$ 的稳定子, 乘法映射在自同构群之间诱导同构. 因而分类映射拟有限. $V$ 的数值有效性由定理~\bookref{6.9.21} 建立. 由 Gabber 引理 (\bookref{5.6.1}), 可选概形的有限覆盖 $Z\to\overline{\mathcal M}_g$. 在 $Z$ 上应用 Kollár 判据 (\bookref{6.9.2}), 得到 $\lambda_{18g-18}=\det Q$ 在 $Z$ 上的拉回丰沛. 因此 $L$ 沿有限态射 $Z\to\overline{\mathcal M}_g\to\overline M_g$ 的拉回丰沛, 从而 $L$ 在 $\overline M_g$ 上丰沛 (见习题~\bookref{5.5.23}).
\end{proof}
\begin{remark}\label{src:6.9.4}
事实上, 稳定曲线 $C$ 的三重典范嵌入 $|\omega_C^{\otimes3}|:C\hookrightarrow\mathbb P^{5g-6}$ 射影正规, 且由二次方程定义; 见~\cite[58 页]{Mum70b}. 因而乘法映射
\[
\operatorname{Sym}^2\pi_*(\omega_{\mathcal U_g/\overline{\mathcal M}_g}^{\otimes3})\twoheadrightarrow\pi_*(\omega_{\mathcal U_g/\overline{\mathcal M}_g}^{\otimes6})
\]
满射, 同样的论证表明 $\lambda_6$ 的充分可除张量幂下降为 $\overline M_g$ 上的丰沛线丛.
\end{remark}
\begin{corollary}[\texorpdfstring{$\overline M_{g,\allowbreak n}$}{带标记稳定曲线粗模空间}的射影性]\label{src:6.9.5}
对满足 $2g-2+n>0$ 的整数 $g,\allowbreak n$, $\overline M_{g,\allowbreak n}$ 在 $\operatorname{Spec}\mathbb Z$ 上射影.
\end{corollary}
\begin{proof}
已知 $\overline M_{0,\allowbreak 3}=\operatorname{Spec}\mathbb Z$, $\overline M_{1,\allowbreak 1}=\mathbb P^1$. 由上述定理, 当 $g\geq2$ 时 $\overline M_g$ 也射影. 因 $\overline{\mathcal M}_{g,\allowbreak n+1}\to\overline{\mathcal M}_{g,\allowbreak n}$ 是万有族 (命题~\bookref{6.6.12}), 它是稳定曲线族, 且 $L:=\omega_{\overline{\mathcal M}_{g,\allowbreak n+1}/\overline{\mathcal M}_{g,\allowbreak n}}$ 相对于 $\overline{\mathcal M}_{g,\allowbreak n}$ 丰沛 (命题~\bookref{6.3.23}). 因而 $L$ 的某个幂下降为 $\overline M_{g,\allowbreak n+1}$ 上相对于 $\overline M_{g,\allowbreak n}$ 丰沛的线丛 (习题~\bookref{5.4.29}(2)). 于是 $\overline M_{g,\allowbreak n+1}\to\overline M_{g,\allowbreak n}$ 为射影态射, 归纳得到当 $2g-2+n>0$ 时 $\overline M_{g,\allowbreak n}$ 射影.
\end{proof}
\begin{remark}\label{src:6.9.6}
另一种证明 $\overline M_{g,\allowbreak n}$ 射影的方法是, 先证明 $\pi_*(L^{\otimes3})$ 数值有效, 再应用 Kollár 判据 (\bookref{6.9.2}); 其中 $L:=\omega_{\mathcal U_{g,\allowbreak n}/\overline{\mathcal M}_{g,\allowbreak n}}(\sigma_1+\cdots+\sigma_n)$, $\sigma_1,\allowbreak \ldots,\allowbreak \sigma_n$ 是万有族的截面; 见~\cite[命题 4.7]{Kol90}.
\end{remark}
近年来, Kollár 判据已用于越来越一般的情形, 例如 Hassett 的加权带标记曲线模空间~\cite{Has03}, 任意维数的稳定簇的模空间~\cite{KP17}, 以及 K-多重稳定 Fano 簇的模空间~\cite{CP21a,XZ20a}.

\noindent\textbf{本节纲要.} 在~\bookref{6.9.2-subsection} 节讨论代数空间的丰沛性判据, 在~\bookref{6.9.3-subsection} 节证明 Kollár 丰沛性判据, 在~\bookref{6.9.4-subsection} 节建立多重典范丛的数值有效性.
\subsection{正性与丰沛性判据}\label{src:6.9.2-subsection}
将~\bookref{B.2.2-subsection} 节引入的线丛正性概念由概形推广到代数空间, 并将~\bookref{B.2.4-subsection} 节的数值有效向量丛概念推广到代数空间.
\begin{definition}\label{src:6.9.7}
拟紧代数叠 $\mathcal X$ 上的线丛 $L$ 称为:
\begin{enumerate}[(1)]
\item \newterm{无基点}{base point free}的, 若对每个 $x\in|\mathcal X|$, 存在 $s\in\Gamma(\mathcal X,\allowbreak L)$ 使 $s(x)\ne0$;
\item \newterm{半丰沛}{semiample}的, 若某个 $n>0$ 使 $L^{\otimes n}$ 无基点.
\end{enumerate}
域 $k$ 上固有代数空间 $X$ 上的线丛 $L$ 称为:
\begin{enumerate}[(1)]\setcounter{enumi}{2}
\item 丰沛的, 若 $X$ 是概形, 且 $L$ 在通常意义下丰沛 (见命题~\bookref{B.2.8});
\item 数值有效的 (分别地, 严格数值有效的), 若对每条整闭曲线\footnote{由定理~\bookref{5.5.32}, 每条这样的整闭曲线都是射影概形.} $C\subseteq X$, 有 $L\cdot C=\deg L|_C\geq0$ (分别地, $>0$);
\item 大的, 若存在常数 $C$, 对所有 $m\geq0$ 均有 $\mathrm h^0(Z,\allowbreak L|_Z^{\otimes m})\geq Cm^{\dim Z}$.
\end{enumerate}
固有代数空间 $X$ 上的向量丛 $E$ 称为数值有效的, 若 $\mathbb P(E)$ 上的 $\mathcal O_{\mathbb P(E)}(1)$ 数值有效.
\end{definition}
若 $f:X'\to\allowbreak X$ 是代数空间的有限满射, 则 $X$ 上的线丛 $L$ 丰沛当且仅当 $f^*L$ 丰沛 (习题~\bookref{5.5.23}).
\begin{lemma}\label{src:6.9.8}
设 $X$ 为域 $k$ 上的固有代数空间.
\begin{enumerate}[(1)]
\item 若 $f:X'\to\allowbreak X$ 是代数空间的固有满射, 则向量丛 $E$ 数值有效当且仅当 $f^*E$ 数值有效.
\item 若向量丛 $E$ 数值有效, 则 $\det E$ 也数值有效.
\item 若 $f:X'\to\allowbreak X$ 是整代数空间的泛拟有限固有态射, 则线丛 $L$ 为大线丛当且仅当 $f^*L$ 为大线丛.
\end{enumerate}
\end{lemma}
\begin{proof}
(1) 只需对线丛 $L$ 证明. 对像为 $C\subseteq X$ 的整子曲线 $C'\subseteq X'$, 由等式 $\deg(f^*L)|_{C'}=\deg(C'\to\allowbreak C)\deg L|_C$ 即得. 对 (2), 应用 Gabber 引理 (\bookref{5.6.1}), 得到概形的有限覆盖 $f:X'\to\allowbreak X$. 由 (1), 只需证明 $f^*\det E=\det(f^*E)$ 数值有效, 而这由概形的情形 (命题~\bookref{B.2.15}) 得到. (3) 可用与概形情形 (命题~\bookref{B.2.24}) 相同的论证证明.
\end{proof}
下题推广引理~\bookref{B.2.26}; 虽然这里不应用它, 它能帮助理解我们的丰沛性证明策略.
\begin{exercise}[容易]\label{src:6.9.9}
设 $\mathcal X$ 为域 $k$ 上的固有 Deligne--Mumford 叠, 粗模空间为 $\mathcal X\to\allowbreak X$. 若 $\mathcal X$ 上的线丛 $L$ 半丰沛且严格数值有效, 则某个 $N>0$ 使 $L^{\otimes N}$ 下降为 $X$ 上的丰沛线丛, 且 $X$ 射影.
\end{exercise}
上述习题的半丰沛性假设在实际中可能很难验证. GIT 方法将半丰沛性直接纳入半稳定性定义 (定义~\bookref{8.2.1}): 给定约化群 $G$ 在射影空间 $\mathbb P(V)$ 上的线性作用, $0\ne v\in V$ 半稳定当且仅当存在 $d>0$ 及 $s\in\Gamma([\mathbb P(V)/G],\allowbreak \mathcal O(d))$ 使 $s(v)\ne0$. 换言之, 半稳定轨迹 $\mathbb P(V)^{\mathrm{ss}}$ 是最大的开子概形, 使 $\mathcal O(1)$ 在 $[(\mathbb P(V)\setminus\mathbb P(V)^{\mathrm{ss}})/G]$ 上半丰沛; 即 $[\mathbb P(V)^{\mathrm{ss}}/G]$ 是 $[\mathbb P(V)/G]$ 上 $\mathcal O(1)$ 的稳定基点轨迹. Hilbert--Mumford 判据 (\bookref{8.4.4}) 提供了验证半稳定性的有效方法.

概形 $X$ 的 Nakai--Moishezon 丰沛性判据 (\bookref{B.2.28}) 从线丛 $L=\mathcal O_X(D)$ 的数值有效性, 及对每个整闭子概形 $Z$ 存在 $L|_Z^{\otimes N}$ 的截面, 得到半丰沛性. 它对代数空间也成立.
\begin{theorem}[Nakai--Moishezon 丰沛性判据]\label{src:6.9.10}
设 $X$ 为代数闭域 $k$ 上的固有代数空间, $L$ 为 $X$ 上的线丛. 以下条件等价:
\begin{enumerate}[(1)]
\item $L$ 丰沛;
\item $L$ 数值有效, 且对每个整闭子概形 $Z\subseteq X$, $L|_Z$ 为大线丛;
\item $L$ 严格数值有效, 且对每个整闭子概形 $Z\subseteq X$, 存在 $m>0$ 使 $L|_Z^{\otimes m}$ 有效.
\end{enumerate}
\end{theorem}
\begin{proof}
与概形情形 (定理~\bookref{B.2.28}) 相同的论证证明非平凡的方向 (3) $\Rightarrow$ (1).
\end{proof}
\begin{remark}[经典表述]\label{src:6.9.11}
不难将交数的定义推广到代数空间. 即证明 $\chi(Z,\allowbreak L|_Z^{\otimes m})$ 是有理多项式, 并将 $c_1(L)^{\dim Z}\cdot Z$ 定义为 $m^{\dim Z}$ 的归一化系数; 见~\cite[标签 0DN3]{SP}. 对数值有效线丛 $L$, 当 $m\gg0$ 时有 $\chi(Z,\allowbreak L|_Z^{\otimes m})=\mathrm h^0(Z,\allowbreak L|_Z^{\otimes m})$, 因而 $L$ 为大线丛等价于 $c_1(L)^{\dim Z}\cdot Z>0$. 这给出代数空间 $X$ 的 Nakai--Moishezon 判据的经典表述:
\[
L\text{ 丰沛}\quad\Longleftrightarrow\quad c_1(L)^{\dim Z}\cdot Z>0\text{ 对每个整闭子空间 }Z\subseteq X.
\]
Seshadri 判据 (\bookref{B.2.30}) 对固有代数空间成立~\cite{Cor93}, 而 Kleiman 判据 (\bookref{B.2.29}) 对一般固有代数空间 (甚至一般固有概形) 尚不清楚.
\end{remark}
\subsection{Kollár 丰沛性判据的证明}\label{src:6.9.3-subsection}
我们讨论并证明 Kollár 判据 (\bookref{6.9.2}): 若 $X$ 为域 $k$ 上的固有代数空间, $W=\operatorname{Sym}^mV\twoheadrightarrow Q$ 为向量丛满射, $w=\operatorname{rk}W$, $q=\operatorname{rk}Q$, $v=\operatorname{rk}V$, 分类映射
\[
X\longrightarrow[\operatorname{Gr}(q,\operatorname{Sym}^mk^v)/\operatorname{PGL}_v],\qquad x\longmapsto[W\otimes\kappa(x)\twoheadrightarrow Q\otimes\kappa(x)]
\]
拟有限, 且 $V$ 数值有效, 则 $\det Q$ 丰沛.
\begin{remark}[更一般的表述]\label{src:6.9.12}
给出另一种表述, 它不仅更一般, 也提示证明的思路. 设 $W\twoheadrightarrow Q$ 为 $X$ 上秩分别为 $w,\allowbreak q$ 的向量丛满射, 满足:
\begin{enumerate}[(a)]
\item $W$ 具有沿仿射代数群 $G$ 的同态 $G\to\operatorname{GL}_w$ 的结构群约化, 且 $G$ 的像在 $w\times w$ 矩阵的射影空间 $\mathbb P(\operatorname{Mat}_{w,\allowbreak w})$ 中的闭包正规;
\item $W$ 数值有效;
\item 分类映射 $X\to[\operatorname{Gr}(q,\allowbreak W)/G]$ 拟有限;
\item 或者 $G$ 线性约化 (如 $\operatorname{char}k=0$ 且 $G$ 约化), 或者 $G=\operatorname{GL}_{v_1}\times\cdots\times\operatorname{GL}_{v_d}$, 且 $G\to\operatorname{GL}_w$ 为半正表示, 即延拓为映射 $\operatorname{Mat}_{v_1,\allowbreak v_1}\times\cdots\times\operatorname{Mat}_{v_d,\allowbreak v_d}\to\operatorname{Mat}_{w,\allowbreak w}$.
\end{enumerate}
则 $\det Q$ 丰沛~\cite[命题 3.6, 引理 3.9]{Kol90}. Kollár 判据从 $\operatorname{Gr}(q,\allowbreak W)$ 上 Plücker 线丛的丰沛性推出 $\det Q$ 丰沛, 推广了在 $\operatorname{Gr}(q,\allowbreak W)$ 上拟有限的固有概形射影这一事实. 的确, 若 $W$ 为平凡向量丛, 结构群可约化为平凡群 $G=\{1\}$, 分类映射 $X\to\operatorname{Gr}(q,\allowbreak W)$ 拟有限便意味着有限, 从而 Plücker 线丛的拉回 $\det Q$ 丰沛. 证明的核心思想是平凡化 $W$, 实质上归约到这一情形.
\end{remark}
\begin{remark}[展开拟有限性条件]\label{src:6.9.13}
由定义~\bookref{4.3.34}, 分类映射 $X\to[\operatorname{Gr}(q,\allowbreak W)/G]$ 拟有限等价于: (i) 对每个代数闭域 $\bar k$, $X(\bar k)\to\operatorname{Gr}(q,\allowbreak W)(\bar k)$ 具有有限纤维; (ii) 对 $x\in X(\bar k)$, 只有有限多个 $G(\bar k)$ 的元素保持 $\ker(\bar k^w\cong W\otimes\bar k\to\allowbreak Q\otimes\bar k)$ 不变.

事实上, 如~\cite[定理 5.1]{KP17} 所观察的, (c) 中分类映射的拟有限性可替换为 $X(k)\to\operatorname{Gr}(q,\allowbreak W)(k)$ 在集合意义下的拟有限性. 令 $\Gamma:X\to\allowbreak X\times[\operatorname{Gr}(q,\allowbreak W)/G]$ 为分类映射的图, 则集合意义下的拟有限性等价于 $\operatorname{im}\Gamma\to[\operatorname{Gr}(q,\allowbreak W)/G]$ 拟有限, 而这足以使下面的证明成立.
\end{remark}
\begin{remark}[稳定性]\label{src:6.9.14}
上述定理不要求 $X$ 的像落在 $\operatorname{Gr}(q,\allowbreak W)$ 的 $G$-稳定轨迹内. 不过, 若确实如此, 则复合 $X\to[\operatorname{Gr}(q,\allowbreak W)^{\mathrm s}/G]\to\operatorname{Gr}(q,\allowbreak W)\mathbin{/\!/}G$ 是固有代数空间的拟有限态射, 因而有限. 所以 $\det Q^{\otimes w}\otimes\det W^{\otimes(-q)}$ 丰沛, 因为它是 $\operatorname{Gr}(q,\allowbreak W)\mathbin{/\!/}G$ 上丰沛线丛的拉回.
\end{remark}
现在证明 Kollár 丰沛性判据~\cite[引理 3.9]{Kol90}. 另见~\cite[定理 5.1]{KP17} 及~\cite{CLM22} 中的精彩阐述. 证明将 $\det Q$ 的丰沛性归约到它的大性; 后者将在命题~\bookref{6.9.15} 中建立.
\begin{proof}[定理~\bookref{6.9.2} 的证明]
由 Nakai--Moishezon 丰沛性判据 (\bookref{6.9.10}), 只需证明 $\det Q$ 数值有效, 且对每个整子空间 $Z\subseteq X$, $\det Q|_Z$ 为大线丛. 因 $W$ 数值有效, 商 $Q$ 也数值有效, 引理~\bookref{6.9.8}(1) 推出 $\det Q$ 数值有效. 若 $Z\subseteq X$ 为整子空间, 则由 Chow 引理 (\bookref{5.6.5}), 存在来自正规射影概形的泛拟有限固有态射 $Z'\to\allowbreak Z$. 考虑在 $Z'$ 上诱导的满射 $W_{Z'}=\operatorname{Sym}^mV_{Z'}\twoheadrightarrow Q_{Z'}$. 分类映射 $Z'\to\allowbreak X\to[\operatorname{Gr}(q,\allowbreak \operatorname{Sym}^mk^v)/\operatorname{PGL}_v]$ 泛拟有限, 且 $V_{Z'}$ 数值有效. 命题~\bookref{6.9.15} 推出 $\det Q_{Z'}$ 为大线丛, 引理~\bookref{6.9.8}(3) 进一步推出 $\det Q$ 为大线丛.
\end{proof}
在 $X$ 射影的假设下, 我们在分类映射泛有限这一较弱假设下建立 $\det Q$ 为大线丛这一较弱结论. 下面的证明利用大性的双有理不变性, 将射影标架丛 $\mathbb P\mathrm{Fr}_W$ 的紧化爆破, 得到到 $\operatorname{Gr}(q,\allowbreak k^w)$ 的态射, 从而构造截面.
\newpage
\begin{proposition}\label{src:6.9.15}
设 $X$ 为域 $k$ 上的正规射影概形. 设 $W=\operatorname{Sym}^mV\twoheadrightarrow Q$ 为向量丛满射, $w=\operatorname{rk}W$, $q=\operatorname{rk}Q$, $v=\operatorname{rk}V$. 假设
\begin{enumerate}[(a)]
\item 分类映射 $X\to[\operatorname{Gr}(q,\allowbreak \operatorname{Sym}^mk^v)/\operatorname{PGL}_v]$ 泛拟有限;
\item $V$ 数值有效.
\end{enumerate}
则 $\det Q$ 为大线丛.
\end{proposition}
\begin{proof}
\textbf{第一步: 平凡化 $V$.} 定义 $X$ 上的相对射影空间
\[
\mathbb P:=\mathbb P(\mathcal Hom_{\mathcal O_X}(V,\mathcal O_X^{\oplus v})),
\]
并令 $\mathbb P\mathrm{Fr}_V\subseteq\mathbb P$ 为自然映射
\begin{equation}\label{src:6.9.16}
\mathcal O_{\mathbb P}^{\oplus v}\longrightarrow V_{\mathbb P}\otimes\mathcal O_{\mathbb P}(1)
\end{equation}
的行列式非零所定义的开子概形. 该映射由万有商 $\mathcal Hom_{\mathcal O_{\mathbb P}}(V_{\mathbb P},\allowbreak \mathcal O_{\mathbb P}^{\oplus w})\to\mathcal O_{\mathbb P}(1)$ 诱导, 其中 $V_{\mathbb P}$ 为 $V$ 沿 $\mathbb P\to\allowbreak X$ 的拉回. 投影 $\mathbb P\mathrm{Fr}_V\to\allowbreak X$ 为主 $\operatorname{PGL}_v$-丛, 事实上正是 $V$ 的射影标架丛 (见习题~\bookref{B.1.57}). 映射~\bookeqref{6.9.16} 在 $\mathbb P\mathrm{Fr}_V$ 上满射, 定义态射 $\mathbb P\mathrm{Fr}_V\to\operatorname{Gr}(q,\allowbreak \operatorname{Sym}^mk^v)$, 置于 Cartesian 图中:
\[
\begin{tikzcd}[column sep=small]
\mathbb P\arrow[dr,"p"']&\mathbb P\mathrm{Fr}_W\arrow[l,hook']\arrow[r]\arrow[d]&\operatorname{Gr}(q,\operatorname{Sym}^mk^v)\arrow[d]\\
&X\arrow[r]\arrow[ur,phantom,"\square" description]&{[\operatorname{Gr}(q,\operatorname{Sym}^mk^v)/\operatorname{PGL}_w]}.
\end{tikzcd}
\]
注意 $\mathbb P\cong\mathbb P((V^\vee)^{\oplus v})$, 故 $p_*\mathcal O_{\mathbb P}(N)$ 等同于 $\operatorname{Sym}^N((V^\vee)^{\oplus v})$. 因 $V$ 数值有效, 数值有效性的性质 (命题~\bookref{B.2.33}(3)) 推出 $(p_*\mathcal O_{\mathbb P}(N))^\vee$ 数值有效. 下面正是在这里使用假设 (b).

\textbf{第二步: 爆破 $\mathbb P$, 将 $\mathbb P\mathrm{Fr}_W\to\operatorname{Gr}(q,\allowbreak \operatorname{Sym}^mk^v)$ 延拓为映射 $\mathbb P'\to\operatorname{Gr}(q,\allowbreak \operatorname{Sym}^mk^v)$.} 若由~\bookeqref{6.9.16} 诱导的复合
\[
\operatorname{Sym}^m\mathcal O_{\mathbb P}^{\oplus v}\longrightarrow\operatorname{Sym}^mV_{\mathbb P}\otimes\mathcal O_{\mathbb P}(m)\longrightarrow Q_{\mathbb P}\otimes\mathcal O_{\mathbb P}(m)
\]
满射, 便不必爆破. 否则取第 $q$ 个外幂, 得到映射 $\bigwedge^q\operatorname{Sym}^m\mathcal O_{\mathbb P}^{\oplus v}\to\det Q_{\mathbb P}\otimes\mathcal O_{\mathbb P}(mq)$. 令 $I$ 为其像子层, 定义 $\pi:\mathbb P'\to\mathbb P$ 为理想层 $I\otimes(\det Q_{\mathbb P})^\vee\otimes\mathcal O_{\mathbb P}(-mq)\subseteq\mathcal O_{\mathbb P}$ 的爆破. 则 $\mathbb P'\to\mathbb P$ 在 $\mathbb P\mathrm{Fr}_W$ 上为同构, 且有映射 $\mathbb P'\to\operatorname{Gr}(q,\allowbreak \operatorname{Sym}^mk^v)$ 延拓 $\mathbb P\mathrm{Fr}_W\to\operatorname{Gr}(q,\allowbreak \operatorname{Sym}^mk^v)$, 得到交换图
\[
\begin{tikzcd}[column sep=small,row sep=large]
E\arrow[d,hook]&H\arrow[d,hook]&L\arrow[d,hook]\\
\mathbb P'\arrow[r,"\pi"']\arrow[rr,bend left=15]&\mathbb P\arrow[r,dashed]\arrow[d,"p"']&\operatorname{Gr}(q,\operatorname{Sym}^mk^v)\arrow[d]\\
&X\arrow[r]&{[\operatorname{Gr}(q,\operatorname{Sym}^mk^v)/\operatorname{PGL}_w]}.
\end{tikzcd}
\]
其中 $E\subseteq\mathbb P'$ 为例外除子, $H\subseteq\mathbb P$ 为超平面类, $L$ 为 $\operatorname{Gr}(q,\allowbreak \operatorname{Sym}^mk^v)$ 上对应于 Plücker 嵌入的丰沛除子. 滥用记号, 也将 $\det Q$ 看作 $X$ 上的除子. 以下标表示除子的拉回, 有公式
\begin{equation}\label{src:6.9.17}
L_{\mathbb P'}=(\det Q)_{\mathbb P'}+mqH_{\mathbb P'}-E.
\end{equation}
\textbf{第三步: 利用泛拟有限性, 对 $X$ 上的丰沛除子 $A$ 及某个 $N>0$, 构造 $\mathcal O_X(m\det Q-A)\otimes\operatorname{Sym}^{Nmq}((V^\vee)^{\oplus v})$ 的截面.} 因 $X\to[\operatorname{Gr}(q,\allowbreak \operatorname{Sym}^mk^v)/\operatorname{PGL}_w]$ 泛拟有限, $\mathbb P'\to\operatorname{Gr}(q,\allowbreak \operatorname{Sym}^mk^v)$ 也泛拟有限. 命题~\bookref{B.2.24} 推出 $L_{\mathbb P'}$ 为大除子. 对每个甚丰沛除子 $A$, 拉回 $A_{\mathbb P'}$ 有效, Kodaira 引理 (\bookref{B.2.20}) 推出存在 $N>0$ 使 $NL_{\mathbb P'}-A_{\mathbb P'}$ 有效. 利用等式~\bookeqref{6.9.17}, 得到线丛的包含
\[
\begin{aligned}
\mathcal O_{\mathbb P'}(NL_{\mathbb P'}-A_{\mathbb P'})
&\cong\mathcal O_{\mathbb P'}(N(\det Q)_{\mathbb P'}+NmqH_{\mathbb P'}-NE-A_{\mathbb P'})\\
&\subseteq\mathcal O_{\mathbb P'}(N(\det Q)_{\mathbb P'}+NmqH_{\mathbb P'}-A_{\mathbb P'})\\
&\cong\pi^*\bigl(p^*\mathcal O_X(N\det Q-A)\otimes\mathcal O_{\mathbb P}(Nmq)\bigr).
\end{aligned}
\]
选取非零截面 $\mathcal O_{\mathbb P'}\to\pi^*(p^*\mathcal O_X(m\det Q-A)\otimes\mathcal O_{\mathbb P}(mq))$. 因 $\mathbb P$ 正规且 $\pi:\mathbb P'\to\mathbb P$ 双有理, $\pi_*\mathcal O_{\mathbb P'}=\mathcal O_{\mathbb P}$. 对所选截面应用 $\pi_*$ 并利用投影公式, 得到非零截面
\[
\mathcal O_{\mathbb P}\longrightarrow p^*\mathcal O_X(N\det Q-A)\otimes\mathcal O_{\mathbb P}(Nmq).
\]
应用 $p_*$, 得到非零截面 $\mathcal O_X\to\mathcal O_X(N\det Q-A)\otimes p_*\mathcal O_{\mathbb P}(Nmq)$, 将它改写为
\begin{equation}\label{src:6.9.18}
(p_*\mathcal O_{\mathbb P}(Nmq))^\vee\longrightarrow\mathcal O_X(N\det Q-A).
\end{equation}
\textbf{第四步: 爆破 $X$ 使~\bookeqref{6.9.18} 满射, 并利用 $V$ 的数值有效性推出 $\det Q$ 为大线丛.} 如前所述, 由命题~\bookref{B.2.33}(3), $V$ 数值有效推出 $(p_*\mathcal O_{\mathbb P}(Nmq))^\vee\cong\operatorname{Sym}^{Nmq}((V^\vee)^{\oplus v})^\vee$ 数值有效. 因而若~\bookeqref{6.9.18} 满射, 则 $N\det Q-A$ 数值有效, 于是 $N\det Q$ 表为数值有效线丛与丰沛线丛之和, 甚至推出 $\det Q$ 丰沛. 否则, 令 $I$ 为~\bookeqref{6.9.18} 的像子层, 爆破理想层 $I\otimes\mathcal O_X(N\det Q-A)^\vee\subseteq\mathcal O_X$, 得到双有理态射 $f:X'\to\allowbreak X$, 例外除子 $E'$, 以及满射
\[
\begin{aligned}
f^*(p_*\mathcal O_{\mathbb P}(Nmq))^\vee&\twoheadrightarrow M,\\
M&:=f^*(\mathcal O_X(N\det Q-A))\otimes\mathcal O_{X'}(-E').
\end{aligned}
\]
因为左方数值有效, $M$ 也数值有效. 这将 $N\det Q_{X'}$ 表为数值有效除子 $M$, 大除子 $A_{X'}$, 有效除子 $E'$ 之和. 由大性的刻画 (\bookref{B.2.21}), $A_{X'}$ 的某个倍数是丰沛除子与有效除子之和. 因而 $\det Q_{X'}$ 的某个倍数是数值有效, 丰沛, 有效除子之和. 数值有效加丰沛仍丰沛, 丰沛加有效为大, 故 $\det Q_{X'}$ 为大线丛. 又因 $X'\to\allowbreak X$ 双有理, $\det Q$ 为大线丛.
\end{proof}
\subsection{多重典范丛的数值有效性}\label{src:6.9.4-subsection}
我们证明 $\overline{\mathcal M}_g$ 上多重典范丛 $\pi_*(\omega_{\mathcal U_g/\overline{\mathcal M}_g}^{\otimes k})$ 的数值有效性: 证明存在负次数线丛商会与 Ekedahl--Mumford 消灭定理矛盾.
\begin{theorem}[Ekedahl--Mumford 消灭定理]\label{src:6.9.19}
设 $X$ 为代数闭域 $k$ 上的极小一般型光滑射影曲面. 若 $n\geq1$, 则 $\mathrm h^1(X,\allowbreak K_X^{\otimes(-n)})\leq1$. 若 $n>1$ 或 $\operatorname{char}k\ne2$, 则 $\mathrm h^1(X,\allowbreak K_X^{\otimes(-n)})=0$.
\end{theorem}
\begin{proof}
特征 $0$ 的版本是~\cite[定理 2]{Mum67}, 正特征情形是~\cite{Eke88} 的主定理.
\end{proof}
\begin{corollary}\label{src:6.9.20}
设 $X$ 为代数闭域 $k$ 上的极小一般型光滑射影曲面. 设 $D$ 为约化有效 Cartier 除子, 每个连通分支的亏格均至少为 $2$. 若 $n\geq2$, 则 $\mathrm h^1(X,\allowbreak K_X^{\otimes n}(D))\leq1$. 若 $n>2$ 或 $\operatorname{char}k\ne2$, 则 $\mathrm h^1(X,\allowbreak K_X^{\otimes n}(D))=0$.
\end{corollary}
\begin{proof}
由短正合列
\[
0\longrightarrow K_X^{\otimes n}\longrightarrow K_X^{\otimes n}(D)\longrightarrow K_X^{\otimes n}(D)|_D\longrightarrow0
\]
及定理~\bookref{6.9.19} 中的不等式, 只需证明当 $n\geq2$ 时 $\mathrm h^1(X,\allowbreak K_X^{\otimes n}(D)|_D)=0$. 由伴随公式, $K_X^{\otimes n}(D)|_D\cong\omega_D^{\otimes n}$. 因 $D$ 每个连通分支的亏格至少为 $2$, 有 $\mathrm h^1(X,\allowbreak K_X^{\otimes n}(D)|_D)=\mathrm h^1(D,\allowbreak \omega_D^{\otimes n})=0$.
\end{proof}
当 $k\geq2$ 时 $\pi_*(\omega_{\mathcal U_g/\overline{\mathcal M}_g}^{\otimes k})$ 的数值有效性归约为光滑曲线上的稳定曲线族的下述性质.
\begin{theorem}\label{src:6.9.21}
若 $T$ 为域 $k$ 上的光滑连通射影曲线, $\pi:C\to\allowbreak T$ 为稳定曲线族, 则当 $k\geq2$ 时 $\pi_*(\omega_{C/T}^{\otimes k})$ 数值有效.
\end{theorem}
\begin{proof}
\textbf{第一步: 归约到特征 $p$.} 设 $\operatorname{char}k=0$. 因 $T,\allowbreak C$ 在 $k$ 上有限型, 其定义方程只涉及 $k$ 的有限多个系数. 因而存在有限生成 $\mathbb Z$-子代数 $A\subseteq k$ 及 Cartesian 图
\[
\begin{tikzcd}
C\arrow[r]\arrow[d]&\widetilde C\arrow[d]\\
T\arrow[r]\arrow[d]\arrow[ur,phantom,"\square" description]&\widetilde T\arrow[d]\\
\operatorname{Spec}k\arrow[r]\arrow[ur,phantom,"\square" description]&\operatorname{Spec}A.
\end{tikzcd}
\]
其中 $\widetilde C,\allowbreak \widetilde T$ 为 $A$ 上有限型概形. 由极限方法 (\bookref{B.3.7} 与~\bookref{6.2.25}), 还可使 $\widetilde T\to\operatorname{Spec}A$ 为光滑曲线族, $\widetilde C\to\widetilde T$ 为稳定曲线族. 最后, 沿映射 $\operatorname{Spec}R\to\operatorname{Spec}A$ 限制, 可设 $A$ 为离散赋值环, 闭点与泛点的特征分别为 $p>0$ 与 $0$. 因向量丛的数值有效性在泛化下稳定 (\bookref{B.2.34}), 只需证明 $\operatorname{char}k=p>0$ 的情形.

\textbf{第二步: 进一步归约.} 我们断言还可假设:
\begin{enumerate}[(a)]
\item $\pi:C\to\allowbreak T$ 泛光滑;
\item $T$ 的亏格至少为 $2$;
\item $C$ 为极小一般型光滑曲面.
\end{enumerate}
归约到 (a) 由习题~\bookref{6.9.23} 处理. 对 (b), 取有限覆盖 $g:T'\to\allowbreak T$, 其中 $T'$ 为亏格 $g\geq2$ 的光滑连通曲线, 令 $C':=C\times_TT'$. 由对偶化层的性质,
\[
g^*\pi_*(\omega_{C/T}^{\otimes k})\cong\pi'_*(\omega_{C'/T'}^{\otimes k}).
\]
由引理~\bookref{6.9.8}(1), $\omega_{C'/T'}^{\otimes k}$ 数值有效推出 $\omega_{C/T}^{\otimes k}$ 数值有效. 对 (c), 若 $f:\widetilde C\to\allowbreak C$ 为奇点解消, 则 $f_*(\omega_{\widetilde C/T}^{\otimes k})\cong\omega_{C/T}^{\otimes k}$, 从而可设 $C$ 光滑. 若 $C$ 含光滑有理 $(-1)$-曲线 $E$, 则 $E$ 必在 $C\to\allowbreak T$ 的某条纤维中, 否则存在有限覆盖 $E=\mathbb P^1\to\allowbreak T$, 而 $T$ 的亏格 $g\geq2$. 由 Castelnuovo 收缩定理 (\bookref{B.2.6}), 有收缩 $E$ 的态射 $f:C\to\allowbreak C'$. 因 $f_*\omega_{C/T}^{\otimes k}=\omega_{C'/T}^{\otimes k}$, 且收缩光滑有理 $(-1)$-曲线的过程有限步后终止, 可设 $C$ 极小. 因 $T$ 与 $C\to\allowbreak T$ 的泛纤维均光滑且为一般型, $C$ 也是一般型.

\textbf{第三步: 正特征情形.} 若 $\pi_*(\omega_{C/T}^{\otimes k})$ 不数值有效, 则存在商线丛 $\pi_*(\omega_{C/T}^{\otimes k})\twoheadrightarrow M$, 其中 $d=\deg M<0$. 考虑绝对 Frobenius 态射 $F:C\to\allowbreak C$, $F:T\to\allowbreak T$, 它们置于交换图
\[
\begin{tikzcd}
C\arrow[r,"F"]\arrow[d,"\pi"']&C\arrow[d,"\pi"]\\
T\arrow[r,"F"']&T.
\end{tikzcd}
\]
注意 $F^*\pi_*(\omega_{C/T}^{\otimes k})=\pi_*(\omega_{C/T}^{\otimes k})$. 因 $\deg F^*M=pd$, 可反复应用 Frobenius, 使 $d$ 任意小. 具体地, 可使 $M^\vee\cong\omega_T^{\otimes k}\otimes L$, 其中 $L$ 为 $T$ 上甚丰沛线丛. (归约到特征 $p$ 的全部目的, 正是反复应用 Frobenius 任意降低 $M$ 的次数.)

满射 $\pi_*(\omega_{C/T}^{\otimes k})\twoheadrightarrow M\cong(\omega_T^{\otimes k}\otimes L)^\vee$ 给出 $T$ 上的向量丛满射
\[
\pi_*(\omega_{C/T}^{\otimes k})\otimes\omega_T^{\otimes k}\otimes L\twoheadrightarrow\mathcal O_T.
\]
因 $\dim T=1$, 得到
\begin{equation}\label{src:6.9.22}
\mathrm h^1(T,\pi_*(\omega_{C/T}^{\otimes k})\otimes\omega_T^{\otimes k}\otimes L)\geq\mathrm h^1(T,\mathcal O_T)\geq2.
\end{equation}
为得到矛盾, 考察 $\mathrm h^1(C,\allowbreak \omega_C^{\otimes k}\otimes\pi^*L)$. 由 Leray 谱序列在 $E_2$ 页退化, 它以 $\mathrm h^1(T,\allowbreak \pi_*(\omega_C^{\otimes k}\otimes\pi^*L))$ 为下界. 得到不等式
\[
\begin{aligned}
\mathrm h^1(C,\omega_C^{\otimes k}\otimes\pi^*L)
&\geq\mathrm h^1(T,\pi_*(\omega_C^{\otimes k}\otimes\pi^*L))\\
&=\mathrm h^1(T,\pi_*(\omega_{C/T}^{\otimes k}\otimes\pi^*\omega_T^{\otimes k}\otimes\pi^*L))\\
&=\mathrm h^1(T,\pi_*(\omega_{C/T}^{\otimes k})\otimes\omega_T^{\otimes k}\otimes L)\\
&\geq\mathrm h^1(T,\mathcal O_T)\geq2.
\end{aligned}
\]
第一个等号用 $\omega_C\cong\omega_{C/T}\otimes\pi^*\omega_T$, 第二个等号用投影公式, 最后用~\bookeqref{6.9.22}. 因 $L$ 甚丰沛, 存在由 $\pi:C\to\allowbreak T$ 的光滑纤维的并组成的有效除子 $D$, 使 $\mathcal O_C(D)\cong\pi^*L$. 上述不等式与将推论~\bookref{6.9.20} 应用于 $D$ 所得结论矛盾. 另见~\cite[定理 4.3]{Kol90} 与~\cite[定理 6.10]{CLM22}.
\end{proof}
\begin{exercise}[细节]\label{src:6.9.23}
设 $T$ 为域 $k$ 上的光滑连通射影曲线. 证明: 若每个泛光滑稳定族的 $\pi_*(\omega_{C/T}^{\otimes k})$ 均数值有效, 则每个稳定族的该向量丛均数值有效.
\end{exercise}
